\clearpage
\begingroup
\fancyhead[L]{\small THEORY OF THE MOTION OF THE MOON.}
\renewcommand{\headrulewidth}{0pt}
\SourceLeaf{626}
\section*{THIRD SECTION.}
\begin{center}
\itshape First approximation.
\end{center}

\par\noindent\hspace*{\parindent}Our first task here is to determine \(m''\).
From the theory of elliptic motion it is known that, if we regard the
motion of the Earth about the Sun solely as a result of the attractive forces of these
two celestial bodies, and set the mean motion of the Sun \(=nu\) and
[the] mean distance of the Earth from the Sun \(=B\), the following relation holds.
\[
m'+m''=nnB^3\frac{\mathrm{d}u^2}{\mathrm{d}t^2};
\]
\noindent Thus, setting \(\dfrac{m'}{m''}=\mu\), a ratio known with sufficient
accuracy \(\left(=\dfrac{1}{365412}\right)\), we obtain
\[
m''=nnB^3\frac{\mathrm{d}u^2}{\mathrm{d}t^2}\frac{1}{1+\mu}
\mathpunct{\raisebox{.4ex}{.}}
\]
\noindent Furthermore,
\[
R'=\BareRoot{\left(RR-\frac{2R}{p}\cos(v-V)+\frac{1+\theta\theta}{pp}\right)}
\]
\noindent or, omitting higher powers of \(\dfrac{1}{R}\) and \(\theta\),
\[
\begin{aligned}
\frac{1}{R'^3}&=\frac{1}{R^3}\left\{1+\frac{3\cos(v-V)}{Rp}\right.\\
&\qquad\left.+\frac{15}{4}\frac{1}{RRpp}(1+\cos2(v-V))
-\frac32\frac{1}{RRpp}\right\}\mathpunct{\raisebox{.4ex}{.}}
\end{aligned}
\]
\noindent This gives
\[
\begin{aligned}
\pi&=\frac{nn}{1+\mu}\frac{B^3}{R^3}\left\{\frac32\sin2(v-V)\right.\\
&\qquad\left.+\frac{1}{Rp}\left(\frac38\sin(v-V)+\frac{15}{8}\sin3(v-V)\right)\right\}\\[5pt]
\omega&=\frac{nn}{1+\mu}\frac{B^3}{R^3}\left\{\frac32+\frac32\cos2(v-V)\right.\\
&\qquad\left.+\frac{1}{Rp}\left(\frac{33}{8}\cos(v-V)+\frac{15}{8}\cos3(v-V)\right)\right\}\\[5pt]
\psi&=\frac{nn}{1+\mu}\frac{B^3}{R^3}\left\{\frac12+\frac32\cos2(v-V)\right.\\
&\qquad\left.+\frac{1}{Rp}\left(\frac98\cos(v-V)+\frac{15}{8}\cos3(v-V)\right)\right\}
\mathpunct{\raisebox{.4ex}{.}}
\end{aligned}
\]
\SourceLeaf{627}
\noindent Finally, if we set the eccentricity of the Sun \(=\varepsilon\) and the mean
anomaly \(=a\), we have
\[
\frac{R}{B}=1+\tfrac12\varepsilon\varepsilon+\varepsilon\cos a
-\tfrac12\varepsilon\varepsilon\cos2a,
\]
\noindent hence
\[
\frac{B^3}{R^3}=1+\tfrac32\varepsilon\varepsilon-3\varepsilon\cos a
+\tfrac92\varepsilon\varepsilon\cos2a,
\]
\noindent and
\[
V=nu-2\varepsilon\sin a+\tfrac54\varepsilon\varepsilon\sin2a.
\]

\par\noindent\hspace*{\parindent}We further note that when substituting elliptic motion
in the terms \(\sin\lambda\) and \(\cos M\), we must keep the symbols \(\lambda\), \(M\)
undetermined. The substitution of \(v-\mathit{Const.}\)
for them is impermissible. The perturbation by the Sun might well impart
to the node and the apogee a motion with a nonperiodic component,
so that \(\lambda\) and \(M\) would take the form \(\zeta v-\mathit{Const.}\),
where \(\zeta\) would differ little from unity.
In that case the assumption of elliptic motion would be altogether false
if we changed \(\lambda\) from \(=\zeta v-\mathit{Const.}\) to \(=v-\mathit{Const.}\),
because this difference would accumulate.
We therefore leave the arguments \(\lambda\) and \(M\) undetermined. More precisely,
we set \(\lambda=\zeta v-\mathit{Const.}\), \(M=\eta v-\mathit{Const.}\),
so that \(\zeta\), \(\eta\) denote coefficients whose value in the absence of the perturbation
by the Sun would be \(=1\) and now remains undetermined
until the calculation settles it.

\par\noindent\hspace*{\parindent}If in this first approximation we set aside the
terms in the third equation that contain \(\dfrac{nn}{R}\), \(nn\varepsilon\) and \(nne\),
the equation becomes
\[
\begin{aligned}
0&=\frac{\mathrm{dd}\theta}{\mathrm{d}v^2}+(1+\tfrac32nn)\theta\\
&\qquad+\tfrac32nn\left(\theta\cos2(v-V)
-\frac{\mathrm{d}\theta}{\mathrm{d}v}\sin2(v-V)\right);
\end{aligned}
\]
\noindent Here we may put for \(V\) the expression \(nv[+\mathit{Const.}]\).
Thus, writing for \(v-nv[-\mathit{Const.}]\) the symbol \(E\) and for
\(\theta\) the expression \(g\sin\lambda\), we obtain
\[
0=\frac{\mathrm{dd}\theta}{\mathrm{d}v^2}+(1+\tfrac32nn)\theta
-\tfrac32nng\sin(2E-\lambda).
\]

\par\noindent\hspace*{\parindent}This equation immediately shows that the expression
\(\lambda=v-\mathit{Const.}\) must be replaced by
\(=v\BareRoot{(1+\tfrac32nn)}-\mathit{Const.}\).
Thus \(\dfrac{\mathrm{d}\lambda}{\mathrm{d}v}=1+\tfrac34nn\), or the
node has a retrograde motion whose ratio to the mean motion
of the Moon is \(\tfrac34nn:1\). From this follows the first approximate
\SourceLeaf{628}value
\[
\theta=g\sin\lambda+(\tfrac38n+\tfrac{3}{32}nn)g\sin(2E-\lambda).
\]

\par\noindent\hspace*{\parindent}To see what numerical coefficients this first
approximation gives, we note that according to observation
\[
\begin{array}{r@{\;}c@{\;}r}
\log n&=&8{,}873\,9096\\
\log nn&=&7{,}747\,8192\\
\log\tfrac34&=&9{,}875\,0613\\\cline{2-3}
&&7{,}622\,8805
\end{array}
\]
\noindent so that \(\tfrac34nn=0{,}004\,1964\,35\), while the observations give
\(\dfrac{\mathrm{d}\lambda}{\mathrm{d}v}=1{,}004\,0218\,706\).

\par\noindent\hspace*{\parindent}Furthermore, \(g\) in arcseconds is \(=18567\overset{\prime\prime}{,}933\).
\[
\begin{array}{r@{\;}c@{\;}r@{\qquad}r@{\;}c@{\;}r}
\log&=&4{,}268\,7636&&&4{,}268\,7636\\
\log\tfrac38n&=&8{,}447\,9409&\log\tfrac{3}{32}nn&=&6{,}719\,7905\\\cline{2-3}\cline{5-6}
&&2{,}716\,7045&&&0{,}988\,5541\\[3pt]
&&520\overset{\prime\prime}{,}840&&&+9\overset{\prime\prime}{,}740
\end{array}
\]
\noindent The coefficient is therefore \(=530\overset{\prime\prime}{,}580\), which, after a
reduction to be explained below, agrees exactly with the observations.
\clearpage
\endgroup
